\documentclass[11pt]{article}
\usepackage[margin=1in]{geometry}
\usepackage{amsmath,amssymb}
\setlength{\parindent}{0pt}
\setlength{\parskip}{0.7em}
\begin{document}
\textbf{21-127 Concepts of Mathematics (fictional)}

Homework 1 - Foundations of proof --- Practice submission A

Student: Fictional learner

\section*{Question 1: Sets in three forms}

(a) $A=\{16,25,27,36,49,64,81\}$.

(b) $B=\{5k+1:k\in\mathbb{Z},\ k\geq0\}$.

(c) $C=\{3^k:k\in\mathbb{Z}\}$.

\newpage
\textbf{21-127 Concepts of Mathematics (fictional)}

Homework 1 - Foundations of proof --- Practice submission A

Student: Fictional learner

\section*{Question 2: A divisibility claim}

False: take $a=2$ and $b=5$. Their product is $10=10\cdot1$. Neither $2/10$ nor $5/10$ is an integer, so neither factor is divisible by $10$. This disproves the universal implication.

\newpage
\textbf{21-127 Concepts of Mathematics (fictional)}

Homework 1 - Foundations of proof --- Practice submission A

Student: Fictional learner

\section*{Question 3: A rational point inside an interval}

My first thought was to call this a weighted average, but I will check the required definition directly. Choose integers $a,b,c,d$ with $b,d\neq0$ such that $x=a/b$ and $y=c/d$.

Then $u=(3ad+bc)/(4bd)$. Both parts are integers, and $4bd$ is nonzero. Hence $u\in\mathbb{Q}$.

Finally, $u-x=(y-x)/4>0$ and $y-u=3(y-x)/4>0$, since $x<y$. This proves the strict inequalities.

\newpage
\textbf{21-127 Concepts of Mathematics (fictional)}

Homework 1 - Foundations of proof --- Practice submission A

Student: Fictional learner

\section*{Question 4: An odd expression}

For any integer $n$, first factor $n^2+3n=n(n+3)$. If $n$ is even, this product is even. If $n$ is odd, write $n=2k+1$ with $k\in\mathbb{Z}$; then $n+3=2(k+2)$, so the product is again even.

Thus $n(n+3)=2m$ for some integer $m$, and $n^2+3n+1=2m+1$. This is odd by definition. No positivity assumption was made.

\newpage
\textbf{21-127 Concepts of Mathematics (fictional)}

Homework 1 - Foundations of proof --- Practice submission A

Student: Fictional learner

\section*{Question 5: Removing either set}

Let $t\in A\setminus(B\cup C)$. It lies in $A$ but in neither $B$ nor $C$, so it lies in both differences on the right. This proves left-to-right containment.

Now let $t$ belong to the intersection on the right. Membership in its two factors says $t\in A$, $t\notin B$, and $t\notin C$. Thus it lies in $A$ but outside $B\cup C$, which is the left-hand side. Both containments hold, so the sets are equal.

\end{document}
